\documentclass[11pt]{amsart}

\usepackage{amsmath}
\usepackage{amssymb}
\usepackage{amsthm}
\usepackage{booktabs}
\usepackage{array}
\usepackage{hyperref}

\newtheorem{theorem}{Theorem}[section]
\newtheorem{proposition}[theorem]{Proposition}
\newtheorem{lemma}[theorem]{Lemma}
\newtheorem{corollary}[theorem]{Corollary}
\theoremstyle{definition}
\newtheorem{definition}[theorem]{Definition}
\theoremstyle{remark}
\newtheorem{remark}[theorem]{Remark}

\numberwithin{equation}{section}
\setlength{\emergencystretch}{2em}

\newcommand{\NN}{\mathbb{N}}
\newcommand{\ZZ}{\mathbb{Z}}
\newcommand{\QQ}{\mathbb{Q}}
\newcommand{\RR}{\mathbb{R}}
\newcommand{\CC}{\mathbb{C}}
\newcommand{\cO}{\mathcal{O}}
\newcommand{\cD}{\mathcal{D}}
\newcommand{\cC}{\mathcal{C}}
\newcommand{\Nm}{\operatorname{N}}
\newcommand{\eu}{\phi_{\mathrm E}}
\newcommand{\Res}{\operatorname{res}}
\newcommand{\lean}[1]{{\def\_{\textunderscore\allowbreak}\texttt{#1}}}

\begin{document}

\hypersetup{pdftitle={The sum of the reciprocal cubes of the Fibonacci numbers is irrational},pdfauthor={Shawn Krisman}}
\title[Reciprocal cubes of the Fibonacci numbers]{The sum of the reciprocal cubes of the Fibonacci numbers is
irrational}

\author{Shawn Krisman}

\subjclass[2020]{Primary 11J72; Secondary 11B39, 33D15, 41A21, 68V20}
\keywords{Fibonacci zeta function, irrationality, $q$-zeta values, Hankel determinants, Lean}

\begin{abstract}
We prove that $\zeta_F(3)=\sum_{n\ge1}F_n^{-3}$ is irrational, where $F_n$ are the Fibonacci numbers. To our
knowledge the only odd value of the Fibonacci zeta function previously known to be irrational was $\zeta_F(1)$. The
proof is a Hankel-determinant construction over $\ZZ[\varphi]$, $\varphi=(1+\sqrt5)/2$, adapted from a recent
unrefereed construction for the $q$-zeta value $\sum n^2/(p^n-1)$ at integers $p\ge2$. Four estimates are combined: a
real decay of size $\exp(-\tfrac29K^3\log\varphi)$, a smallness $\exp(-K^3\log\varphi)$ at the Galois conjugate, which
replaces the base-prime content gain of the integer case, local cyclotomic denominator bounds at the twisted roots of
unity that occur, and nonvanishing by rational interpolation. The whole argument is formalised in Lean~4 with
Mathlib.
\end{abstract}

\maketitle

% ===========================================================================
\section{Introduction}
% ===========================================================================

Let $F_n$ be the Fibonacci numbers and $\zeta_F(s)=\sum_{n\ge1}F_n^{-s}$ the Fibonacci zeta function. Andr\'e-Jeannin
proved that $\zeta_F(1)$ is irrational \cite{AJ89}, by Pad\'e approximation of a Lambert series. Duverney,
Ke.~Nishioka, Ku.~Nishioka and Shiokawa proved that the even values $\zeta_F(2s)$ are transcendental, and so are
$\sum_n F_{2n-1}^{-s}$ for every $s\ge1$ \cite{DNNS97}; Elsner, Shimomura and Shiokawa determined the algebraic
relations among these numbers \cite{ESS07,ESS08}. All of these sums are values of quasi-modular forms, to which
Nesterenko's theorem \cite{Nes96} applies. The odd values $\zeta_F(2s+1)$, $s\ge1$, are not of this kind, and as far
as we know nothing was proved about their arithmetic (see the survey \cite{Wal10} and the introduction of
\cite{MS24}).

\begin{theorem}\label{thm:main}
$\displaystyle\zeta_F(3)=\sum_{n\ge1}\frac1{F_n^3}$ is irrational.
\end{theorem}

With $t=\varphi^{-1}=(\sqrt5-1)/2$ and Binet's formula, $F_n=t^{-n}(1-(-t^2)^n)/\sqrt5$, so
\begin{equation}\label{eq:Xi}
\Xi:=\frac{\zeta_F(3)}{5\sqrt5}=\sum_{N\ge1}\frac{t^{3N}}{(1-(-t^2)^N)^3}.
\end{equation}
This is a $q$-series at the quadratic unit $t$, a relative of the $q$-zeta value
$\zeta_q(3)=\sum_n n^2q^n/(1-q^n)$ \cite{KRZ06}. For $q=1/p$ with $p\ge2$ an integer, the irrationality of
$\zeta_q(3)$ remains open in the refereed literature: Krattenthaler, Rivoal and Zudilin proved that one of
$\zeta_q(3),\dots,\zeta_q(11)$ is irrational \cite{KRZ06}, Jouhet and Mosaki reduced this to
$\zeta_q(3),\dots,\zeta_q(9)$ \cite{JM10}, and Fischler and Zudilin proved linear independence results of the same
kind \cite{FZ10}. In October 2026 a proof that $\sum_{n\ge1}n^2/(p^n-1)$ is irrational for every integer $p\ge2$ was
claimed in an unrefereed, AI-assisted research draft with a Lean formalisation, posted on GitHub under the pen name
uplift-universal-inc \cite{UPL,QZ3}. Its construction is a moving-denominator Hankel determinant in the line of
Uvarov \cite{Uva69}, Krattenthaler \cite{Kra21}, Pr\'evost \cite{Pre96}, Zudilin \cite{Zud16,Zud17} and the unrefereed
preprint of Fauzan \cite{Fau}.

Expanding \eqref{eq:Xi} in Lambert series does not reduce Theorem~\ref{thm:main} to a statement about $\zeta_q(3)$: one
obtains a $\QQ(\sqrt5)$-linear combination of $1$, odd parts of $\zeta_q(3)$ and $\zeta_q(1)$ at $q=\varphi^{-2}$, and
values of holomorphic modular forms of weights $1$ and $3$ at $q=\varphi^{-1},\varphi^{-2}$. We therefore treat the
series \eqref{eq:Xi} directly.

\subsection*{The method}
We adapt the construction of \cite{UPL} to the real quadratic field $\QQ(\sqrt5)$. A linear functional $L$, a second
Euler derivative summed against the weights $(-t)^n$ at the nodes $(-t^2)^n$, has rational moments and pure simple-pole
tails of $\Xi$ (Lemma~\ref{lem:moments}). The Hankel determinants $\Delta_K(X)$ of $L$ against the moving denominators
$D_K(s)=\prod_{a\le K}(1-(-t^2)^as)$ are polynomials of degree $\le K$ in $X$ with $\Delta_K(\Xi)=\det[L(s^{i+j}/D_K)]$.
Four estimates are combined:
\begin{center}
\small
\begin{tabular}{lll}
\toprule
 & size (per $K^3\log\varphi$) & \\
\midrule
real decay $|\Delta_K(\Xi)|$ & $-2/9$ & Proposition~\ref{prop:decay}\\
conjugate $|\sigma\Delta_K(-\Xi)|$ & $-1$ & Proposition~\ref{prop:conj}\\
denominators $\log|\Nm C_K(t)|$ & $\le c_*=1.1954\ldots<11/9$ & Propositions~\ref{prop:global}, \ref{prop:cost}\\
nonvanishing & infinitely many $K$ & Proposition~\ref{prop:nonvan}\\
\bottomrule
\end{tabular}
\end{center}
If $\zeta_F(3)$ were rational, $C_K(t)(25v)^K\Delta_K(\Xi)$ would be a nonzero element of $\ZZ[\varphi]$ of norm
$<1$ for large $K$.

Three ingredients are new compared with \cite{UPL}. First, the setting: the Fibonacci kernel is a single $q$-series over
$\ZZ[\varphi]$, handled with a formal parameter $\tau$ specialised at $\tau=t$ and at its conjugate $\tau=-\varphi$. Second,
since $\varphi$ is a unit there is no base prime and no coefficient content to extract. Its role is taken by the
smallness of the conjugate determinant at the conjugate point $\sigma\Xi=-\Xi$ (Proposition~\ref{prop:conj}). Its size,
$K^3\log\varphi=\tfrac12K^3\log\varphi^2$, is the size of the content gain $\tfrac12K^3\log p$ of \cite{UPL} at
$p=\varphi^2$. Third, the poles of the coefficients sit at two kinds of twisted roots of unity, those $\zeta$ for which
$-\zeta^2$ has small order and those for which $-\zeta(-\zeta^2)^k=1$; we prove local denominator bounds of both types
(Section~\ref{sec:den}). The present paper does not depend logically on \cite{UPL}: every step is proved here, and the
formal proof is self-contained.

\subsection*{Formal verification}
The theorem and every step of its proof are formalised in Lean~4 with Mathlib
(\url{https://github.com/telaviv/fibonacci-zeta3-irrationality}); Section~\ref{sec:lean} gives the
statement, the map from the propositions below to Lean declarations, and what has to be trusted. The formal proof uses
the elementary denominator constant $c_*=\frac{11}8\beta_7+\frac13=\frac{4393}{3675}$ of Proposition~\ref{prop:cost},
where $\beta_7=\prod_{p\le7}(1-p^{-2})$; it needs no information on the distribution of Euler's function beyond a
sieve by $2,3,5,7$.

% ===========================================================================
\section{The functional and the determinants}\label{sec:setup}
% ===========================================================================

\subsection*{Notation}
Throughout, $\varphi=(1+\sqrt5)/2$ is the golden ratio, $t=\varphi^{-1}$, $\cO=\ZZ[\varphi]$, and $\sigma$ is the
nontrivial automorphism of $\QQ(\sqrt5)$, so $\sigma t=-\varphi$. Euler's function is written $\eu(d)$ and the
M\"obius function $\mu(d)$. We put
\[
\rho=-t,\qquad w=-t^2,\qquad \alpha_k=\rho w^k\ (k\ge0),
\]
\[
\psi_N=\frac{t^{3N}}{(1-w^N)^3},\qquad S_a=\sum_{N=1}^a\psi_N,
\]
$A_K=K(K+1)/2$, and $\theta=s\,\frac{d}{ds}$. Then $\rho^2=-w$, $|w|=t^2$ and $\Xi=\sum_N\psi_N$. We write $h_m$ for the
complete homogeneous symmetric polynomial of degree $m$ ($h_m=0$ for $m<0$). Constants in $O(\cdot)$ are absolute and
$K\ge2$.

\begin{definition}
For $g$ analytic on a neighbourhood of $|s|\le t^2$ put
$\displaystyle L(g)=\sum_{n\ge1}\rho^n\,[\theta(\theta+1)g](w^n)$.
\end{definition}
The series converges absolutely since $\theta(\theta+1)g=O(s)$ and $|\rho^nw^n|=t^{3n}$.

\begin{lemma}\label{lem:moments}
\textup{(i)} $L(s^k)=\nu_k:=k(k+1)\alpha_k/(1-\alpha_k)$ for $k\ge0$.
\textup{(ii)} For $a\ge1$, $L\bigl(1/(1-w^as)\bigr)=\lambda_a(\Xi-S_a)$ with $\lambda_a=2\rho^{-a}$.
\end{lemma}
\begin{proof}
(i) $\theta(\theta+1)s^k=k(k+1)s^k$; sum the geometric series. (ii) $\theta(\theta+1)(1-cs)^{-1}=2cs/(1-cs)^3$, so the
left side is $\sum_n\rho^n\,2w^{a+n}/(1-w^{a+n})^3=2\rho^{-a}\sum_{N>a}\rho^Nw^N/(1-w^N)^3$, and
$\rho^Nw^N=t^{3N}$.
\end{proof}

Let $D_K(s)=\prod_{a=1}^K(1-w^as)$, with poles $x_a=w^{-a}$, $|x_a|=\varphi^{2a}$. For a polynomial $P$ of degree
$<2K-1$,
\[
\frac{P}{D_K}=Q+\sum_{a=1}^K\frac{\Res_a(P)}{1-w^as},\qquad \Res_a(P)=\frac{P(x_a)}{\prod_{b\ne a}(1-w^{b-a})},
\]
with $Q$ a polynomial of degree $\le\deg P-K$. Let $L_X$ be the linear map with $L_X(s^k)=\nu_k$ and
$L_X(1/(1-w^as))=\lambda_a(X-S_a)$, $X$ an indeterminate, and
\[
\Delta_K(X)=\det\bigl[L_X(s^{i+j}/D_K)\bigr]_{0\le i,j<K}\in\QQ(\sqrt5)[X],\qquad\deg\Delta_K\le K.
\]
By Lemma~\ref{lem:moments}, $\Delta_K(\Xi)=\det[L(s^{i+j}/D_K)]$.

\subsection*{The formal parameter}
Every quantity above is the value at $\tau=t$ of a rational function of an indeterminate $\tau$ with rational
coefficients: replace $t$ by $\tau$, so $\rho=-\tau$, $w=-\tau^2$, $\alpha_k=(-1)^{k+1}\tau^{2k+1}$ and
$\psi_N=\tau^{3N}/(1-(-\tau^2)^N)^3$. Write $\Delta_{K,\tau}(X)\in\QQ(\tau)[X]$ for the resulting polynomial, and
$w(\zeta)=-\zeta^2$, $\alpha_k(\zeta)$ for the values at $\tau=\zeta$. Since $\sigma t=-\varphi$, the coefficientwise
conjugate $\sigma\Delta_K$ is the value of $\Delta_{K,\tau}$ at $\tau=-\varphi$. Both values are defined because no
denominator vanishes at a real number that is neither $0$ nor a root of unity.

% ===========================================================================
\section{Real decay}\label{sec:decay}
% ===========================================================================

\begin{proposition}\label{prop:decay}
$\log|\Delta_K(\Xi)|\le-\tfrac29(K^3+3K^2-4)\log\varphi+O(K^2)$.
\end{proposition}
\begin{proof}
Let $V_j(s)=\prod_{u=1}^j(s-w^{\lceil u/3\rceil})$, $R_j=\lceil j/3\rceil$ and $a_j=\sum_{u\le j}\lceil u/3\rceil$.
The change of basis from $1,s,\dots,s^{K-1}$ to $V_0,\dots,V_{K-1}$ is unitriangular, so
$\Delta_K(\Xi)=\det[L(V_iV_j/D_K)]$.

Since $\theta(\theta+1)g(s_0)$ depends only on the $2$-jet of $g$ at $s_0$, the $m$-th term of $L(V_iV_jW)$,
$W=1/D_K$, vanishes for $m<R:=\max(R_i,R_j)$: then $V_iV_j$ has a triple zero at $w^m$. For $m\ge R$ and $s=w^m$ every
root $\beta$ of $V_j$ satisfies $|s|\le|\beta|$, so $|s-\beta|\le2|\beta|$ and $\theta(s-\beta)=s$. Hence
$|V_j(s)|\le2^j|w|^{a_j}$ and $|\theta^uV_j(s)|\le j^u2^j|w|^{a_j}|s|/|w|^{R_j}$ for $u=1,2$. On $|s|\le t^2$ we have
$|w^as|\le t^{2a+2}$, so $|W|$, $|\theta W|/|s|$ and $|\theta^2W|/|s|$ are bounded by an absolute constant. Since
$\theta(\theta+1)=\theta^2+\theta$ kills constants, by Leibniz every term of $\theta(\theta+1)(V_iV_jW)$ contains at
least one $\theta$ on some factor, hence a factor $\le|s|/|w|^R$. With $|\rho|^m\le1$ and
$\sum_{m\ge R}|w|^{m-R}\le(1-t^2)^{-1}$,
\[
|L(V_iV_jW)|\le c_1(i+j+1)^24^{i+j}|w|^{a_i+a_j}
\]
for an absolute constant $c_1$. Pull $|w|^{a_i}$ out of row $i$ and column $i$ and apply Hadamard's inequality:
$|\Delta_K(\Xi)|\le|w|^{2\sum_{j<K}a_j}(c_1(2K)^24^{2K}\sqrt K)^K$. Finally
$2\sum_{j<K}a_j=(K^3+3K^2-\iota_K)/9$ with $\iota_K\in\{0,4,2\}$ according to $K\bmod3$.
\end{proof}

% ===========================================================================
\section{Denominators}\label{sec:den}
% ===========================================================================

\subsection{The integrality ring and the root types}

\begin{lemma}\label{lem:ring}
The coefficients of $\Delta_{K,\tau}(X)$ lie in $\ZZ[\tau,\tau^{-1},\Phi_d(\tau)^{-1}:d\in\cD_K]$, where $\cD_K$ is
the set of orders of roots of unity $\zeta$ such that
\textup{(W)} $w(\zeta)$ has order $e\le K$, or \textup{($\alpha$)} $\alpha_k(\zeta)=1$ for some $1\le k\le K-2$.
\end{lemma}
\begin{proof}
The inputs are the Laurent monomials $w^a$ and $\lambda_a=2(-\tau)^{-a}$; the factors $1/(1-w^{b-a})$ and
$(1-w^N)^{-3}$ ($N\le K$), where $1-w^n=1-(-1)^n\tau^{2n}$ is $\pm$ a product of cyclotomic polynomials $\Phi_d(\tau)$
with $w(\zeta)^n=1$; and $\nu_k=k(k+1)\alpha_k/(1-\alpha_k)$, of which only $k\le K-2$ occur (the quotient has degree
$\le 2K-2-K$), with $1-\alpha_k=1\mp\tau^{2k+1}=\pm\prod\Phi_d(\tau)$. Division by $D_K$ is integral because its leading
coefficient $\prod_a(-w^a)=\pm\tau^{K(K+1)}$ is a unit. No rational constant is inverted.
\end{proof}

For $d\in\cD_K$ the type depends only on $d$. If $\zeta$ has order $d$, then $w(\zeta)$ has order $e(d)=2d$ ($d$ odd),
$d$ ($d\equiv2\bmod4$), $d/4$ ($d\equiv4\bmod8$), $d/2$ ($d\equiv0\bmod8$), and for each $e$ the orders $d$ with
$e(d)=e$ satisfy $\sum_{e(d)=e}\eu(d)=2\eu(e)$. The least $k\ge1$ with $\alpha_k(\zeta)=1$ is
$k_{\min}(d)=(d-1)/2$ ($d\equiv3\bmod4$), $(3d-1)/2$ ($d\equiv1\bmod4$), $(d-2)/4$ ($d\equiv2\bmod8$, $d\ge10$),
$(3d-2)/4$ ($d\equiv6\bmod8$), and $k_{\min}(2)=2$ (the formula $(d-2)/4$ would give $0$); there is none if $4\mid d$.
(Every primitive $d$-th root is a power $\zeta^j$ with $j$ coprime to $2d$; then $-\zeta^{2j}=(-\zeta^2)^j$ and
$\alpha_k(\zeta^j)=\alpha_k(\zeta)^j$, and powers coprime to the order preserve the order.) All $d\in\cD_K$ satisfy
$d\le4K$: $w(\zeta)^e=1$ gives $\zeta^{4e}=1$, and $\alpha_k(\zeta)=1$ gives $\zeta^{4k+2}=1$.

\emph{Local setting.} Fix $\zeta$ of order $d\in\cD_K$, let $\varpi=\tau-\zeta$ and let $v=\operatorname{ord}_\varpi$ on
$\QQ(\zeta)((\varpi))$, extended to polynomials in $X$ as the Gauss valuation (the minimum over the coefficients; it is
multiplicative). Then:
\begin{itemize}
\item[(a)] if $e(d)=e\le K$, then $x_a(\zeta)=x_b(\zeta)$ iff $e\mid a-b$, and then $v(x_a-x_b)=1$: indeed
$x_a-x_b=x_a(1-w^{b-a})$, and $1-w^n$ has only simple zeros, since $\frac{d}{d\tau}w^n\ne0$ at $\tau\ne0$;
\item[(b)] $\psi_N$ has a pole of order exactly $3$ if $e\mid N$ and is regular otherwise, by (a); so $S_a$ has a pole of
order $\le3$, and is regular if $a<e$;
\item[(c)] $1-\alpha_k=1\mp\tau^{2k+1}$ is separable, so it has at most a simple zero at $\zeta$;
\item[(d)] $\lambda_a$ and $x_a$ are Laurent monomials, hence units.
\end{itemize}

\subsection{Local bounds}

Put $c(m)=2\sum_{j=0}^{m-2}\max(0,\tfrac{m+2}2-j)$ and $c'(m)=2\sum_{j=0}^{m-1}\max(0,\tfrac{m+2}2-j)$; thus
$c(1)=0$, $c(m)=4(m-1)$ for $m\le6$, $c(m)\le(m+3)^2/4$, $c'(1)=3$ and $0\le c'-c\le3$.

\begin{proposition}\label{prop:W}
Let $e=e(d)\le K$ and partition $\{1,\dots,K\}$ into the clusters $I_r=\{a\equiv r\bmod e\}$ of sizes $m_r$,
$1\le r\le e$. Then
\[
-v(\Delta_{K,\tau})\le\sum_{r=1}^{e-1}c(m_r)+c'(m_e)+\epsilon\max(0,K-1-k_{\min}(d)),
\]
with $\epsilon=1$ if $\zeta$ is also an $\alpha$-root and $\epsilon=0$ otherwise.
\end{proposition}
\begin{proof}
Let $F(s)=\prod_a(s-x_a)$ and $F_r=\prod_{a\in I_r}(s-x_a)$. Fix a cluster $r$, write $m=m_r$ and
$I_r=\{a_1<\dots<a_m\}$.

\emph{Newton cluster basis.} For $0\le j<m$ put $N_{r,j}=(F/F_r)\prod_{l=m-j+1}^m(s-x_{a_l})$; it vanishes at the top
$j$ nodes of the cluster. At $\varpi=0$ these reduce to $(F/F_r)(s-y_r)^j\cdot$(unit), $y_r=x_r(\zeta)$, which form a
basis of the polynomials of degree $<K$ by the Chinese remainder theorem (distinct residues $r$ give distinct points
$y_r$). So the change of basis from $1,\dots,s^{K-1}$ has unit determinant, and $v(\Delta_{K,\tau})$ is the valuation of
$G=[L_X(NN'/D_K)]$ in this basis.

\emph{Entries.} $G_{(r,j),(r',k)}=P+\sum_a\Res_a(NN')\lambda_a(X-S_a)$, where $P$ is $L_X$ of the polynomial part.
(i) If $r\ne r'$ then $F\mid N_{r,j}N_{r',k}$ and all $\Res_a=0$.
(ii) If $r=r'$, then $\Res_{a_i}=0$ unless $i\le m-\max(j,k)$, and for such $i$
\[
\Res_{a_i}=\frac{(F/F_r)(x_{a_i})^2\prod_{l>m-j}(x_{a_i}-x_{a_l})\prod_{l>m-k}(x_{a_i}-x_{a_l})}
{\prod_{b\ne a_i}(1-w^{b-a_i})}.
\]
Here $(F/F_r)(x_{a_i})$ is a unit, each displayed difference has valuation exactly $1$ by (a), and the denominator has
valuation exactly $m-1$. So $v(\Res_{a_i})=j+k-(m-1)$.
(iii) $v(\lambda_a)=0$, and by (b) the Gauss valuation of $X-S_a$ is $\ge-3$ if $a\ge e$ and $\ge0$ if $a<e$.
(iv) For $r<e$ the node $a_1=r$ is regular and $a_i$ is singular for $i\ge2$; for $r=e$ every node is singular. So the
proper part of the $(j,k)$ entry has $v\ge j+k-(m-1)-3$ if some singular node has $i\le m-\max(j,k)$, and
$v\ge j+k-(m-1)$ otherwise. The second case occurs exactly when $r<e$ and $\max(j,k)=m-1$.

\emph{Weights.} Let $\kappa_j=\max(0,\tfrac{m+2}2-j)$ for $j\le m-2$ and $\kappa_{m-1}=0$ if $r<e$, and
$\kappa_j=\max(0,\tfrac{m+2}2-j)$ for all $j<m$ if $r=e$. We claim, following the two cases of (iv):
\begin{itemize}
\item[(A)] if the $(j,k)$ entry has a singular node (all pairs when $r=e$; the pairs with $\max(j,k)\le m-2$ when
$r<e$), then $\kappa_j+\kappa_k\ge m+2-j-k$;
\item[(B)] if $r<e$ and $\max(j,k)=m-1$, then $\kappa_j+\kappa_k\ge0\ge m-1-j-k$.
\end{itemize}
In case (A) both indices carry the weight $\max(0,\tfrac{m+2}2-j)$ (for $r=e$ this includes
$\kappa_{m-1}=\max(0,\tfrac{4-m}2)$). If both weights are positive the inequality is an equality. If $\kappa_k=0$, then
$k\ge\frac{m+2}2$, so $\kappa_j\ge\frac{m+2}2-j\ge m+2-j-k$. If both vanish, then $j,k\ge\frac{m+2}2$ and the right side
is $\le0$. Case (B) holds because $j+k\ge m-1$.

Scale row $(r,j)$ by $\varpi^{\lceil\kappa_j\rceil}$ and column $(r,j)$ by $\varpi^{\lfloor\kappa_j\rfloor}$. The
required exponents are integers, and $a+b\ge n\in\ZZ$ implies $\lceil a\rceil+\lfloor b\rfloor\ge n$, so by (iv), (A)
and (B) every scaled proper entry is integral. The total scaling is $\sum(\lceil\kappa_j\rceil+\lfloor\kappa_j\rfloor)
=2\sum\kappa_j$, which is $c(m)$ for $r<e$ and $c'(m)$ for $r=e$. All exponents are $\ge0$.

\emph{Polynomial part.} The quotient of $NN'$ by $D_K$ has degree $\le K-2$ and coefficients in $\ZZ[\tau^{\pm1}]$, so
$P=\sum_{k\le K-2}(\text{coeff})\nu_k$. At a pure W-root every $\nu_k$, $k\le K-2$, is regular, so $P$ is integral, and
so is the scaled matrix; this gives the bound with $\epsilon=0$.

At an $\alpha$-root, put $P_\alpha=\{k\le K-2:\alpha_k(\zeta)=1\}$ and write $\nu_k=\varpi^{-1}u_k$ with $u_k$ regular for
$k\in P_\alpha$ (by (c)). Split $P=A^{\mathrm{pol}}+\varpi^{-1}B$, where
$B(f,g)=\sum_{k\in P_\alpha}u_k(0)\,[s^k](fg/D_K)_+$ (here $(\cdot)_+$ is the polynomial part) and $A^{\mathrm{pol}}$
collects $\varpi^{-1}(u_k-u_k(0))[s^k](\cdot)_+$ for $k\in P_\alpha$ and $\nu_k[s^k](\cdot)_+$ for $k\notin P_\alpha$;
both $A^{\mathrm{pol}}$ and $B$ have regular entries. If $\deg f\le k_{\min}$, then $\deg(fg)\le k_{\min}+K-1$, so the
quotient has degree $<k_{\min}$ and $B(f,\cdot)=0$. Since the polynomials of degree $\le k_{\min}$ form a space of
dimension $k_{\min}+1$, $\operatorname{rank}B\le\max(0,K-1-k_{\min})$. After scaling, the matrix is $A+\varpi^{-1}B'$.
Here $A$ is the sum of the scaled proper parts, integral by (A) and (B), and the scaled $A^{\mathrm{pol}}$, integral
because all exponents are $\ge0$. The matrix $B'$ is $B$ in the Newton basis, scaled; its entries are regular and its
rank is that of $B$, since the scalings are invertible over the fraction field. Expand the determinant multilinearly in
the columns: a term taking more than $\max(0,K-1-k_{\min})$ columns from $\varpi^{-1}B'$ vanishes by linear
dependence, and every other term has valuation $\ge-\max(0,K-1-k_{\min})$.
\end{proof}

\begin{lemma}\label{lem:alpha}
If $e(d)>K$ and $\zeta$ is an $\alpha$-root, then $-v(\Delta_{K,\tau})\le\max(0,K-1-k_{\min}(d))$.
\end{lemma}
\begin{proof}
Work in the monomial basis. Since $e>K$, no difference $|a-b|<K$ is divisible by $e$, so by (a)--(d) the coefficients
$\Res_a$, the $\lambda_a$ and the $S_a$ ($a\le K<e$) are regular at $\zeta$. Only the $\nu_k$, $k\in P_\alpha$, carry
poles, and the rank argument of Proposition~\ref{prop:W} applies with no scaling.
\end{proof}

\begin{proposition}\label{prop:global}
Let $b_{K,d}$ be the bound of Proposition~\ref{prop:W} or Lemma~\ref{lem:alpha} \textup{(}$0$ if $d\notin\cD_K$\textup{)}
and $C_K(\tau)=\prod_d\Phi_d(\tau)^{b_{K,d}}$. Then $C_K(\tau)\Delta_{K,\tau}(X)\in\ZZ[\tau^{\pm1}][X]$; hence
$C_K(t)\Delta_K(X)\in\cO[X]$, and $C_K(-\varphi)\sigma\Delta_K(X)\in\cO[X]$ is its conjugate.
\end{proposition}
\begin{proof}
By Lemma~\ref{lem:ring} a coefficient is $P/(\tau^n\prod\Phi_d^{n_d})$ with $P\in\ZZ[\tau]$. The bound $v_\zeta\ge-b_{K,d}$
at every primitive $d$-th root forces $\Phi_d^{n_d-b_{K,d}}\mid P$ in $\QQ[\tau]$, hence in $\ZZ[\tau]$ since $\Phi_d$ is
monic. Specialise $\tau\mapsto t$, a unit of $\cO$, and $\tau\mapsto-\varphi$.
\end{proof}

\subsection{\texorpdfstring{The size of $C_K(t)$}{The size of C\_K(t)}}

\begin{lemma}\label{lem:phid}
For $d\ge1$, $\bigl|\log|\Phi_d(t)\Phi_d(-\varphi)|-\eu(d)\log\varphi\bigr|\le2\mathcal E_1$, where
$\mathcal E_1=\sum_{a\ge1}|\log(1-\varphi^{-a})|=2.1136\ldots$.
\end{lemma}
\begin{proof}
$\Phi_d(x)=\prod_{a\mid d}(x^a-1)^{\mu(d/a)}$ and $\sum_{a\mid d}\mu(d/a)a=\eu(d)$. For $|x|=\varphi$ write
$x^a-1=x^a(1-x^{-a})$, and $|\log|1-x^{-a}||\le|\log(1-\varphi^{-a})|$; for $x=t$ use $|1-t^a|$ in the same way.
\end{proof}

Let $\beta_7=\prod_{p\le7}(1-p^{-2})=\frac{768}{1225}$ and $c_*=\frac{11}8\beta_7+\frac13=\frac{4393}{3675}=1.19537\ldots$.

\begin{proposition}\label{prop:cost}
$\log|\Nm C_K(t)|\le c_*K^3\log\varphi+O(K^2\log K)$.
\end{proposition}
\begin{proof}
By Lemma~\ref{lem:phid}, $\log|\Nm C_K(t)|\le\log\varphi\sum_db_{K,d}\eu(d)+2\mathcal E_1\sum_db_{K,d}$. Here
$\sum_db_{K,d}=O(K^2\log K)$: for each $e\le K$ at most four $d$ have $e(d)=e$, each with a W part
$\le4K+K^2/(4e)+3$ (see below), and the at most $4K$ orders $d\le4K$ have $\alpha$ parts $\le K$. Split
$b_{K,d}$ into its W part and its $\alpha$ part.

\emph{W part.} Since $m_r\ge1$, $\sum_rm_r=K$ and $m_r<K/e+1$, and since $c(m)=4(m-1)+\delta(m)$ with
$0\le\delta(m)\le\max(0,m-5)^2/4$ (for $m\ge6$, $c(m)=(m+2)(m+4)/4$ or $(m+3)^2/4$ according to the parity of $m$),
\[
\sum_{r<e}c(m_r)+c'(m_e)\le4(K-e)+\frac{\max(0,K-4e)^2}{4e}+3.
\]
Summing over the orders $d$ with $e(d)=e$ gives the weight $\sum_{e(d)=e}\eu(d)=2\eu(e)$. Bound $\eu(e)\le e\,\eta(e)$
with $\eta(e)=\prod_{p\le7,\,p\mid e}(1-1/p)$, a function of period $P=210$ with $0\le\eta\le1$ and mean $\beta_7$. So
the W part is at most $\sum_{e\le K}\eta(e)f_K(e)$, where
\[
f_K(y)=8y(K-y)+\tfrac12\max(0,K-4y)^2+6y .
\]
We use the elementary estimate: if $\eta$ has period $P$, $0\le\eta\le1$ and mean $\beta$, and $f$ is a function on
$[0,K]$ with $|f|\le M$ and total variation $V$, then
\[
\Bigl|\sum_{1\le e\le K}\eta(e)f(e)-\beta\int_0^Kf(y)\,dy\Bigr|\le2P(M+V).
\]
(Cut $[0,K]$ into blocks of length $P$. On a full block $(e_0,e_0+P]$ both $\sum\eta(e)f(e)$ and
$\beta\int f$ differ from $P\beta f(e_0)$ by at most $P$ times the oscillation of $f$ on the block; the oscillations
add up to at most $V$, and the last, incomplete block costs at most $2PM$.) For $f_K$ we have $M,V=O(K^2)$, so the W
part is at most
\[
\beta_7\int_0^Kf_K(y)\,dy+O(K^2)=\Bigl(\tfrac43+\tfrac1{24}\Bigr)\beta_7K^3+O(K^2)=\tfrac{11}8\beta_7K^3+O(K^2).
\]

\emph{$\alpha$ part.} $\max(0,K-1-k_{\min})$ counts the $k\in[k_{\min},K-2]$, so
\begin{align*}
\sum_d\max(0,K-1-k_{\min}(d))\,\eu(d)&\le\sum_{k=1}^{K-2}\sum_{j=1}^k\#\{\zeta:\alpha_j(\zeta)=1\}\\
&=\sum_{k=1}^{K-2}(k^2+2k)\le\tfrac13K^3+O(K^2),
\end{align*}
since $\alpha_j(\zeta)=1$ means $\zeta^{2j+1}=(-1)^{j+1}$, which has $2j+1$ solutions.
\end{proof}

\begin{remark}\label{rem:sharp}
The following is a sketch and is not used in the proof. Using $\sum_{n\le y}\eu(n)=3y^2/\pi^2+O(y\log y)$ in residue
classes modulo $8$ and the exact $k_{\min}$, the same bound holds with $c_*$ replaced by
$(12I+80/27)/\pi^2\le1.130790$, where $I=\int_0^1\ell(x)x\,dx\in[0.6822834,0.6831236]$ and $\ell$ is the piecewise
linear function with $\ell(1/m)=c(m)/m$. Numerically, $\sum_db_{K,d}\eu(d)/K^3$ is $1.19$ at $K=20$ and $1.16$ at
$K=40$, decreasing towards $1.1308$.
\end{remark}

% ===========================================================================
\section{Smallness at the conjugate}\label{sec:conj}
% ===========================================================================

\begin{proposition}\label{prop:conj}
$\log|\sigma\Delta_K(-\Xi)|\le-K^2(K+1)\log\varphi+K\log(\sqrt K\,M_K)$, where
$M_K=C_0\bigl(K^2/((1-t^3)(1-t^2))+\varphi^2\bigr)$ and $C_0=\prod_{i\ge1}(1-t^{2i})^{-1}=2.0770906\ldots$.
\end{proposition}
\begin{proof}
At $\tau=-\varphi$ the data become: $w\mapsto\sigma w=w^{-1}$, $\lambda_a\mapsto\sigma\lambda_a=2t^a$,
$S_a\mapsto\sigma S_a=-S_a$ (as $\sigma\psi_N=-\psi_N$) and $\nu_k\mapsto\sigma\nu_k=-k(k+1)/(1+\alpha_k)$ (as
$\sigma\alpha_k=-1/\alpha_k$). Let $L^\sigma_X$ be the corresponding formal functional:
$L^\sigma_X(s^k)=\sigma\nu_k$ and $L^\sigma_X(1/(1-w^{-a}s))=2t^a(X+S_a)$. The conjugate poles are $r_a=w^a$,
$|r_a|=t^{2a}$, and $\sigma D_K(s)=\prod_a(1-w^{-a}s)=(-1)^Kw^{-A_K}F^\sigma(s)$ with $F^\sigma(s)=\prod_a(s-r_a)$. Hence
\[
\sigma\Delta_K(X)=(-1)^{K^2}w^{A_KK}\det[Y_{i+j}(X)],\qquad Y_k(X)=L^\sigma_X(s^k/F^\sigma).
\]
Partial fractions give
$s^k/F^\sigma=\sum_{j=0}^{k-K}h_{k-K-j}(r)s^j+\sum_ar_a^k/\bigl((F^\sigma)'(r_a)(s-r_a)\bigr)$.

\emph{The pole moments at $X=-\Xi$.} Since $1/(1-w^{-a}s)=-r_a/(s-r_a)$, we have
$L^\sigma_X(1/(s-r_a))=-r_a^{-1}\,2t^a(X+S_a)$, and at $X=-\Xi$ this is
$2t^ar_a^{-1}(\Xi-S_a)=2t^ar_a^{-1}\sum_{N>a}\psi_N$. Writing $N=a+M$, $t^{3N}=t^{3a}t^{3M}$, $t^{4a}=r_a^2$ and
$w^{a+M}=r_aw^M$, this equals $\Gamma(r_a)$, where
\[
\Gamma(r)=2r\sum_{M\ge1}\frac{t^{3M}}{(1-rw^M)^3}=\sum_{n\ge1}\gamma_nr^n,\qquad
\gamma_n=n(n+1)\sum_{M\ge1}t^{3M}w^{M(n-1)},
\]
analytic on $|r|<\varphi^2$, with $|\gamma_n|\le n(n+1)t^{2n+1}/(1-t^{2n+1})\le n(n+1)t^{2n+1}/(1-t^3)$. Since
$\sum_ar_a^m/(F^\sigma)'(r_a)=h_{m-K+1}(r)$,
\[
Y_k(-\Xi)=\sum_{j=0}^{k-K}h_{k-K-j}(r)\,\sigma\nu_j+\sum_{n\ge1}\gamma_nh_{k+n-K+1}(r).
\]

\emph{Bounds.} As $h_m(r_1,\dots,r_K)=w^mh_m(1,w,\dots,w^{K-1})$ and the coefficients of
$\prod_{i\ge0}(1-|w|^iu)^{-1}=\sum_m u^m\prod_{i=1}^m(1-t^{2i})^{-1}$ are at most $C_0$, we get $|h_m(r)|\le C_0t^{2m}$.
For $j\ge1$, $|\sigma\nu_j|\le j(j+1)/(1-t^3)$, so the first sum is at most $C_0K^2/((1-t^3)(1-t^2))$ for $k\le2K-2$.
For the second, $\sum_nn(n+1)t^{2n+1}=2t^3/(1-t^2)^3=2$ (as $1-t^2=t$) and $1/(1-t^3)=\varphi^2/2$, so
$\sum_n|\gamma_n|\le\varphi^2$. Hence $|Y_k(-\Xi)|\le M_K$ for $0\le k\le2K-2$. Hadamard's inequality and
$|w|^{A_KK}=\varphi^{-K^2(K+1)}$ finish the proof.
\end{proof}

\begin{remark}
Coefficientwise, the factor $2t^a=2|r_a|^{1/2}$ is not analytic in $r_a$; only at the point $X=\sigma\Xi=-\Xi$ does it
combine with the tail $\sum_{N>a}\psi_N\approx t^{3a}$ into $r_a^2$ times an analytic function. This is why the
conjugate point, rather than the coefficients of $\sigma\Delta_K$, carries the gain.
\end{remark}

% ===========================================================================
\section{Nonvanishing}\label{sec:nonvan}
% ===========================================================================

\begin{proposition}\label{prop:nonvan}
$\Delta_K(\Xi)\ne0$ for infinitely many $K$.
\end{proposition}
\begin{proof}
Let $\cC(z)=L(1/(z-s))=\sum_n\rho^n\,2w^nz/(z-w^n)^3$ for $|z|>t^2$. It continues to $\CC\setminus\{0\}$ with a pole
of order exactly $3$ at each $w^n$ (leading coefficient $2\rho^nw^{2n}\ne0$), so it is not rational; it is analytic at
$\infty$ and $O(1/z)$. For $\deg Q\le n-1$ put $H_{n,Q}(z)=L\bigl(Q(s)/((z-s)D_n(s))\bigr)$, analytic for $|z|>t^2$
and in particular at every $x_a$, and
\[
T_{n,Q}(z)=L\Bigl(\frac{Q(z)D_n(s)-D_n(z)Q(s)}{(z-s)D_n(s)}\Bigr),
\]
a polynomial in $z$ of degree $\le n-1$. Then $\cC Q-T_{n,Q}=D_nH_{n,Q}$.

If $\Delta_n(\Xi)=\det[L(s^{i+j}/D_n)]=0$, a kernel vector gives $Q_n\ne0$ of degree $\le n-1$ with
$L(s^jQ_n/D_n)=0$ for $j<n$. Put $T_n=T_{n,Q_n}$ and $E_n=\cC Q_n-T_n$. For $|z|>t^2+\varepsilon$ the series
$(z-s)^{-1}=\sum_js^jz^{-j-1}$ converges, together with its first two $s$-derivatives, uniformly on $|s|\le t^2$, so $L$
may be applied termwise; hence $H_{n,Q_n}(z)=\sum_{j\ge n}L(s^jQ_n/D_n)z^{-j-1}=O(z^{-n-1})$. So $E_n=O(z^{-1})$ and
$E_n(x_a)=0$ for $a\le n$.

If $\Delta_n(\Xi)=\Delta_{n+1}(\Xi)=0$, then $V=T_{n+1}Q_n-T_nQ_{n+1}=Q_{n+1}E_n-Q_nE_{n+1}$ is a polynomial,
$O(z^{n-1})$, vanishing at $x_1,\dots,x_n$; so $V=0$ and $T_n/Q_n=T_{n+1}/Q_{n+1}$. If $\Delta_n(\Xi)=0$ for all
$n\ge N$, all $T_n/Q_n$ equal one reduced fraction $T/Q$, so $Q_n=U_nQ$ and $T_n=U_nT$, and $E:=\cC Q-T\ne0$ since
$\cC$ is not rational. Then $U_nE=E_n=O(1/z)$ forces $E=c_0z^{-q}(1+o(1))$ with $c_0\ne0$, $q\ge1$, and
$\deg U_n\le q-1$. But $E(x_a)\ne0$ for all large $a$, while $E_n(x_a)=0$; so $U_n(x_a)=0$ for all large $a\le n$, too
many zeros for a polynomial of degree $<q$ once $n$ is large.
\end{proof}

% ===========================================================================
\section{\texorpdfstring{Proof of Theorem~\ref{thm:main}}{Proof of the theorem}}\label{sec:proof}
% ===========================================================================

Suppose $\zeta_F(3)=u/v$ with $u\in\ZZ$, $v\in\ZZ_{>0}$. Since $\sqrt5=1+2t$, \eqref{eq:Xi} gives
$\Xi=u(1+2t)/(25v)$ and $\sigma\Xi=-\Xi$. For $K$ in the infinite set of Proposition~\ref{prop:nonvan} put
$y_K=C_K(t)(25v)^K\Delta_K(\Xi)$. By Proposition~\ref{prop:global}, $y_K\in\cO$, and $y_K\ne0$ because
$\Delta_K(\Xi)\ne0$ and $C_K(t)\ne0$ ($t$ is not a root of unity, so $\Phi_d(t)\ne0$). Moreover
$\sigma y_K=C_K(-\varphi)(25v)^K\sigma\Delta_K(-\Xi)$, so $|\Nm(y_K)|=|y_K\,\sigma y_K|\ge1$, and by
Propositions~\ref{prop:decay}, \ref{prop:conj} and \ref{prop:cost}
\begin{multline*}
0\le\log|\Nm C_K(t)|+2K\log(25v)+\log|\Delta_K(\Xi)|+\log|\sigma\Delta_K(-\Xi)|\\
\le\bigl(c_*-\tfrac29-1\bigr)K^3\log\varphi+O(K^2\log K).
\end{multline*}
Since $c_*-\tfrac{11}9=\tfrac{11}8\beta_7-\tfrac89=-0.02684\ldots<0$, this fails for large $K$. \qed

\begin{remark}
Not used in the proof. With the sharp constant of Remark~\ref{rem:sharp} the margin is
$\tfrac{11}9-1.130790\ge0.0914$ per $K^3\log\varphi$. At finite $K$, computed with exact determinants and interval
arithmetic, the sum $\log|\Nm C_K(t)|+\log|\Delta_K(\Xi)|+\log|\sigma\Delta_K(-\Xi)|$ (without the term
$2K\log(25v)$) is already negative, about $-0.12K^3\log\varphi$ for $8\le K\le26$; and the local bounds of
Section~\ref{sec:den} agree with the exact multiplicities of the cyclotomic factors in $\QQ(\tau)[X]$ for small $K$.
The scripts are available from the author.
\end{remark}

% ===========================================================================
\section{Formal verification}\label{sec:lean}
% ===========================================================================

The Lean development is the repository
\begin{center}
\url{https://github.com/telaviv/fibonacci-zeta3-irrationality};
\end{center}
all file names below refer to it. It (library \lean{FibZeta}, Lean~4.35.0, Mathlib v4.35.0, about 8500 lines) proves
\begin{center}
\lean{theorem FibZeta.irrational\_zetaF3 : Irrational FibZeta.zetaF3},
\end{center}
where \lean{zetaF3} is the series $\sum_{n\ge0}1/F_{n+1}^3$, written with Mathlib's \lean{Nat.fib} and a \lean{tsum}
(\lean{FibZeta/Final.lean}; \lean{hasSum\_zetaF3} shows that the series converges to it). The proof uses only the
axioms \lean{propext}, \lean{Classical.choice} and \lean{Quot.sound} (\lean{Axioms.lean}); there is no \lean{sorry},
no \lean{native\_decide} and no added axiom.

\emph{What has to be trusted.} For the theorem itself, only the definition of \lean{zetaF3} (and Mathlib's
\lean{Nat.fib}, \lean{tsum} and \lean{Irrational}) has to be read: the kernel checks everything else. Since a divergent
\lean{tsum} has the value $0$, which is rational, the statement cannot hold for a degenerate reason. To see that the
formal proof follows the argument of this paper one also reads the definitions in \lean{FibZeta/Defs.lean} and the
intermediate statements in \lean{FibZeta/Statements.lean}; an internal comparison of these with an earlier draft of
this section is in the repository (\lean{FIDELITY\_REFEREE.md}).

\subsection*{Definitions}
\lean{FibZeta/Defs.lean} defines the construction of Section~\ref{sec:setup} over an arbitrary field $\mathbb k$ with a
parameter $\tau\in\mathbb k$: \lean{alpha}, \lean{nu}, \lean{psi}, \lean{S}, the poles \lean{node} $=w^{-a}$, the monic
polynomial \lean{F} $=\prod_a(s-x_a)$ (so $D_K=(-1)^Kw^{A_K}F$), the residues \lean{res}, the two coefficients
\lean{mConst}, \lean{mLin} of $L_X(s^k/F)$, the normalisation \lean{nrm}, and the determinant \lean{Delta}.
Specialised at $\tau=t$ (\lean{tR}) and at $\tau=-\varphi$ it gives $\Delta_K$ and $\sigma\Delta_K$; at
$\tau=$ \lean{RatFunc.X} it gives $\Delta_{K,\tau}$. The real functional is \lean{Lfun}, with
$\theta(\theta+1)g=s^2g''+2sg'$ written with \lean{deriv}, and \lean{XiR} is $\Xi$. The clearing polynomial is
\lean{CKpoly} $=\prod_{d\le4K}\Phi_d^{b_{K,d}}$, with exponents \lean{nExp} $=$ \lean{nW} $+$ \lean{nAlpha}: the W part
(cluster costs \lean{cW}, \lean{cW'} and cluster sizes \lean{clusterSize}) and the $\alpha$ part. The root types are
defined through the chosen root $\exp(2\pi i/d)$ (\lean{zetaC}, \lean{eW}, \lean{IsAlphaRoot}), and
\lean{FibZeta/Roots.lean} proves that they are the same for every primitive $d$-th root.

\subsection*{The chain}
\lean{FibZeta/Statements.lean} states the ingredients as propositions, and \lean{FibZeta.irrational\_zetaF3\_of} (in
\lean{FibZeta/Main.lean}) derives the theorem from them by the argument of Section~\ref{sec:proof}. Instead of norms
in $\cO$ it uses that $t$ and $-\varphi$ are the two roots of $\tau^2+\tau-1$: for $Q\in\ZZ[X]$ the product
$Q(t)Q(-\varphi)$ is an integer, nonzero if $Q(t)\ne0$ (\lean{norm\_ne\_zero}). The docstrings of the Lean files refer
to the numbering of the author's notes; the header of \lean{FibZeta/Statements.lean} translates it to the numbering of
this paper.
\begin{center}
\small
\begin{tabular}{lll}
\toprule
Paper & Lean statement & Lean proof\\
\midrule
Lemma~\ref{lem:moments}, entries & \lean{EntryIdentity} & \lean{FibZeta.Analytic.entryIdentity}\\
Prop.~\ref{prop:decay} & \lean{DecayBound} & \lean{FibZeta.Analytic.decayBound}\\
Lemma~\ref{lem:ring}--Prop.~\ref{prop:global} & \lean{Integrality} & \lean{FibZeta.Denom.integrality}\\
Lemma~\ref{lem:phid}, Prop.~\ref{prop:cost} & \lean{CostBound} & \lean{FibZeta.Cost.costBound}\\
Prop.~\ref{prop:conj} & \lean{ConjBound} & \lean{FibZeta.Conj.conjBound}\\
Prop.~\ref{prop:nonvan} & \lean{Nonvanishing} & \lean{FibZeta.NonVan.nonvanishing}\\
Section~\ref{sec:proof} & & \lean{FibZeta.irrational\_zetaF3\_of}\\
\bottomrule
\end{tabular}
\end{center}
\lean{EntryIdentity} is used only in the proof of \lean{DecayBound}, not by the assembly. The formal statements are
equivalent to, or slightly weaker than, the propositions, in the form the assembly needs. \lean{DecayBound} is
$|\Delta_K(\Xi)|\le C^{K^2}t^{2b_K}$ with $b_K=2\sum_{j<K}a_j$ (proved with $C=4096$). \lean{ConjBound} is
$|\sigma\Delta_K(-\Xi)|\le C^{K^2}t^{K^2(K+1)}$; its proof bounds the complete homogeneous sums $h_m(r)$ by $\varphi^K$
instead of $C_0t^{2m}$, which still gives entries of size $\le D^K$. \lean{CostBound} is
$\log|C_K(t)C_K(-\varphi)|\le\frac{121}{100}K^3\log\varphi+CK^2(1+\log K)^2$; its proof follows
Proposition~\ref{prop:cost}, with cubic constant $c_*=\frac{4393}{3675}$ and the constant $12$ in place of
$2\mathcal E_1$ in the upper half of Lemma~\ref{lem:phid}, and $c_*<\frac{121}{100}<\frac{11}9$. \lean{Integrality} is
Proposition~\ref{prop:global} for every real $\tau\ne0$ that is not a root of unity, with one integer polynomial for
all such $\tau$; this covers $t$ and $-\varphi$ at once. Its proof works in the local rings of $\CC(\tau)$ at the roots
of unity: the Newton cluster basis at W-roots, the monomial basis elsewhere, and the rank correction at
$\alpha$-roots (\lean{FibZeta.Denom.local\_bound\_cluster}, \lean{FibZeta.Denom.local\_bound\_monomial},
\lean{FibZeta.Denom.integrality\_of\_local\_bounds}). The proof of \lean{Nonvanishing} uses the Cauchy transform $\cC$
directly: its values at the interpolation nodes, $\cC(x_a)=2t^a(\Xi-S_a)$, and its expansion
$\cC(1/u)=\sum_k\nu_ku^{k+1}$ at $\infty$ replace the functional $L$ in the interpolation identities.

\subsection*{Provenance}
Parts of the formal development are adapted from the Lean formalisation QZeta3 \cite{QZ3} of \cite{UPL} (Apache
License~2.0): the localisation, specialisation and descent machinery, the weighted and rank determinant lemmas, the
Cauchy-transform and normality arguments for nonvanishing, the decay estimates, the totient sieve and the cyclotomic
size bounds. Every adapted file says so in its header, and the repository's \lean{NOTICE} file records it.

\subsection*{Use of AI tools}
This work was carried out with substantial assistance from AI systems (Anthropic's Claude models), which were used to
search the literature, to develop and check the proof, to write the Lean formalisation and to draft this text. The
author directed the project and is responsible for its content. The formal proof is checked by the Lean kernel.

% ===========================================================================
\bibliographystyle{amsplain}
\bibliography{refs}

\end{document}
